\documentclass[a4paper,twoside]{article}

\usepackage{morewrites}

\usepackage[svgnames,dvipsnames,x11names]{xcolor}
\usepackage{graphicx}
\usepackage[english]{babel}
\usepackage[colorlinks=true, linkcolor=NavyBlue, urlcolor=NavyBlue, citecolor=NavyBlue]{hyperref}
\usepackage[a4paper]{geometry}
\usepackage{ts-fonts}
\usepackage{booktabs}
\usepackage{csquotes}
\usepackage{tabularx}
\usepackage{tcolorbox}
\usepackage{xcolor}
\usepackage{fvextra}
\usepackage{amsmath}
\usepackage{float}
\usepackage{seqsplit}
\newcolumntype{Y}{>{\centering\arraybackslash}X}
\newcolumntype{Z}{>{\raggedleft\arraybackslash}X}
\usepackage{ts-callouts}
\usepackage{ts-code}

\usepackage{lastpage}
\usepackage{refcount}
\makeatletter
\def\ps@plain{
\let\@oddhead\@empty
\let\@evenhead\@empty
\def\@oddfoot{
\hfil
\ifnum\getpagerefnumber{LastPage}>1\relax
\thepage
\fi
\hfil}
\let\@evenfoot\@oddfoot
}
\makeatother

\title{Custom counters}
\author{}\date{}
\begin{document}
\maketitle
Technical documents number things that \LaTeX{} knows nothing about: findings,
requirements, bugs, risks, test cases. TeXSmith lets you declare such a series
in the front matter, mark each item in the body, and reference it anywhere ---
the numbers are computed by TeXSmith rather than by the backend, so the \LaTeX{}
build, the Typst build and --- with the companion plugin --- the MkDocs site all
show the same values.
\begin{tscode}[lang=markdown, engine=pygments]
\begin{Verbatim}[commandchars=\\\{\},breaklines, breakanywhere, commandchars=\\\{\}]
\PYZhy{}\PYZhy{}\PYZhy{}
counters:
  n:
    name: Requirement
\PY{g+gu}{    format: \PYZdq{}N\PYZhy{}\PYZob{}n:02d\PYZcb{}\PYZdq{}}
\PY{g+gu}{\PYZhy{}\PYZhy{}\PYZhy{}}

| Id | Requirement |
| \PYZhy{}\PYZhy{}\PYZhy{} | \PYZhy{}\PYZhy{}\PYZhy{} |
| \PYZsh{}\PYZob{}n:joy\PYZcb{} | Everyone shall be happy |
| \PYZsh{}\PYZob{}n:respect\PYZcb{} | Everyone shall respect the others |

Smile in every circumstance (@n:joy) and do no harm to others (@n:respect).
\end{Verbatim}
\end{tscode}
renders as
\begin{center}
\begin{tabularx}{\linewidth}{>{\raggedright\arraybackslash}X>{\raggedright\arraybackslash}X}
\toprule
\textbf{Id} & \textbf{Requirement} \\
\midrule
N-01 & Everyone shall be happy \\
N-02 & Everyone shall respect the others \\
\bottomrule
\end{tabularx}
\end{center}
Smile in every circumstance (N-01) and do no harm to others (N-02).

Both \tscodeinline{N-\allowbreak{}01} occurrences are the same PDF anchor: the table cell is the target,
the parenthesised one is a clickable link.
\section{Declaring a counter}
Each key under \tscodeinline{counters:} is the \textbf{prefix} used by the markers.
\begin{tscode}[lang=yaml, engine=pygments]
\begin{Verbatim}[commandchars=\\\{\},breaklines, breakanywhere, commandchars=\\\{\}]
\PY{n+nt}{counters}\PY{p}{:}
\PY{+w}{  }\PY{n+nt}{n}\PY{p}{:}
\PY{+w}{    }\PY{n+nt}{name}\PY{p}{:}\PY{+w}{ }\PY{l+lScalar+lScalarPlain}{Requirement}\PY{+w}{       }\PY{c+c1}{\PYZsh{} human\PYZhy{}readable name, used in diagnostics}
\PY{+w}{    }\PY{n+nt}{format}\PY{p}{:}\PY{+w}{ }\PY{l+s}{\PYZdq{}}\PY{l+s}{N\PYZhy{}\PYZob{}n:02d\PYZcb{}}\PY{l+s}{\PYZdq{}}\PY{+w}{     }\PY{c+c1}{\PYZsh{} optional, defaults to \PYZdq{}\PYZob{}n\PYZcb{}\PYZdq{}}
\PY{+w}{    }\PY{n+nt}{start}\PY{p}{:}\PY{+w}{ }\PY{l+lScalar+lScalarPlain}{1}\PY{+w}{                }\PY{c+c1}{\PYZsh{} optional, defaults to 1}
\end{Verbatim}
\end{tscode}
\tscodeinline{format} is a Python format string. Three fields are available:
\begin{center}
\begin{tabularx}{\linewidth}{>{\raggedright\arraybackslash}X>{\raggedright\arraybackslash}X}
\toprule
\textbf{Field} & \textbf{Meaning} \\
\midrule
\tscodeinline{n} & the counter value (\tscodeinline{1}, \tscodeinline{2}, ...), so \tscodeinline{\{n:02d\}} pads it \\
\tscodeinline{prefix} & the counter prefix (\tscodeinline{n}) \\
\tscodeinline{key} & the item key (\tscodeinline{joy}) \\
\bottomrule
\end{tabularx}
\end{center}
The format string is validated when the document is parsed; an invalid one
(\tscodeinline{\{oops\}}, unbalanced braces) raises a \tscodeinline{CounterValidationError}.

Unlike \tscodeinline{glossary:}, a counter has no string shorthand: each entry is a mapping,
because a series without a \tscodeinline{format} is rarely what you want.

A prefix must match \tscodeinline{[A-\allowbreak{}Za-\allowbreak{}z][A-\allowbreak{}Za-\allowbreak{}z0-\allowbreak{}9\_-\allowbreak{}]*} and must not shadow one of the
conventional built-in prefixes --- \tscodeinline{sec}, \tscodeinline{fig}, \tscodeinline{tbl}, \tscodeinline{eq}, \tscodeinline{lst}, \tscodeinline{app},
\tscodeinline{chap}, \tscodeinline{part}, \tscodeinline{note} --- which stay reserved for regular labels.
\section{Marking an item}
\subsection{\tscodeinline{\#\{prefix:key\}} --- define and print}
The marker prints the formatted number and becomes the reference target.
\begin{tscode}[lang=markdown, engine=pygments]
\begin{Verbatim}[commandchars=\\\{\},breaklines, breakanywhere, commandchars=\\\{\}]
\PY{g+gh}{\PYZsh{}\PYZob{}fw:boot\PYZhy{}loop\PYZcb{} The firmware reboots when the watchdog fires.}
\end{Verbatim}
\end{tscode}
It is a plain inline construct: it works in a paragraph, a table cell, a list
item, an admonition title, a heading. Numbers are allocated in document order.

The marker is inert unless its prefix is declared: \tscodeinline{\#\{name\}} (no colon),
\tscodeinline{\#\{sh:var\}} with \tscodeinline{sh} undeclared, or \tscodeinline{\#\{n:joy\}} in a document without a
\tscodeinline{counters:} section all stay literal text. This keeps Ruby/CoffeeScript
interpolations (\tscodeinline{\#\{user.name\}}) intact. Inside code spans and fenced blocks
nothing is ever substituted, and \tscodeinline{\textbackslash{}\#\{n:joy\}} forces a literal.
\subsection{\tscodeinline{\{\#prefix:key\}} --- define silently}
Any element carrying an \tscodeinline{id} whose prefix is declared is numbered too, without
printing anything. This is the usual \tscodeinline{attr\_list} syntax, so it attaches to
headings, figures and tables:
\begin{tscode}[lang=markdown, engine=pygments]
\begin{Verbatim}[commandchars=\\\{\},breaklines, breakanywhere, commandchars=\\\{\}]
\PY{g+gu}{\PYZsh{}\PYZsh{} Boot loop \PYZob{}\PYZsh{}fw:boot\PYZhy{}loop\PYZcb{}}

![\PY{n+nt}{Watchdog trace}](\PY{n+na}{trace.png})\PYZob{}\PYZsh{}fw:trace\PYZcb{}
\end{Verbatim}
\end{tscode}
\tscodeinline{@fw:boot-\allowbreak{}loop} then resolves to \tscodeinline{FW-\allowbreak{}01} even though the number appears nowhere
in the heading.
\begin{tscallout}[kind=warning, title={Position matters for \tscodeinline{\{\#…\}}}]
\tscodeinline{attr\_list} consumes \tscodeinline{\{\#…\}} at the end of a heading, on an image, or on a
line of its own after a block. In the middle of a sentence it stays literal
text --- use \tscodeinline{\#\{…\}} there.
\end{tscallout}
\section{Referencing an item}
\tscodeinline{@prefix:key} (or \tscodeinline{@[prefix:key]}) prints the formatted number as a hyperlink.
This is the regular cross-reference shorthand; a declared
prefix simply routes it to the counter registry.
\begin{tscode}[lang=markdown, engine=pygments]
\begin{Verbatim}[commandchars=\\\{\},breaklines, breakanywhere, commandchars=\\\{\}]
The watchdog issue (@fw:boot\PYZhy{}loop) is fixed in 1.4.2.
\end{Verbatim}
\end{tscode}
The reference prints the number alone --- \tscodeinline{FW-\allowbreak{}01}, not \tscodeinline{Finding FW-\allowbreak{}01}. Write the
noun yourself, as you would for a section. The \tscodeinline{name:} field is only used in
diagnostics.
\section{Diagnostics}
\begin{center}
\begin{tabularx}{\linewidth}{>{\raggedright\arraybackslash}X>{\raggedright\arraybackslash}X}
\toprule
\textbf{Situation} & \textbf{Behaviour} \\
\midrule
\tscodeinline{@n:missing} --- no such item & warning, the reference renders empty \\
\tscodeinline{\#\{n:joy\}} twice with the same key & warning, both print the first number \\
\tscodeinline{\#\{x:joy\}} --- undeclared prefix & left as literal text, no warning \\
invalid \tscodeinline{format} or prefix & \tscodeinline{CounterValidationError} at parse time \\
\bottomrule
\end{tabularx}
\end{center}
\section{Scope and stability}
Numbers are allocated per conversion, in document order, and shared across all
the documents of a multi-document build --- \tscodeinline{a.md}, \tscodeinline{b.md} and \tscodeinline{c.md} continue a
single series rather than restarting.
\begin{tscallout}[kind=danger, title={Numbers are positional}]
Inserting an item renumbers every item after it. When the identifiers leave
the document --- a finding quoted in an audit report, a requirement cited in a
test plan --- that renumbering breaks the external traceability silently.
Pin the values with a dedicated \tscodeinline{start:} per counter and stable ordering, or
keep the volatile numbering for internal documents only. Explicit pinning is
planned but not implemented yet.
\end{tscallout}
\section{Backend mapping}
\begin{center}
\begin{tabularx}{\linewidth}{>{\raggedright\arraybackslash}X>{\raggedright\arraybackslash}X>{\raggedright\arraybackslash}X}
\toprule
\textbf{} & \textbf{Definition} & \textbf{Reference} \\
\midrule
\LaTeX{} & \tscodeinline{\textbackslash{}phantomsection\textbackslash{}label\{n:joy\}N-\allowbreak{}01} & \tscodeinline{\textbackslash{}hyperref[n:joy]\{N-\allowbreak{}01\}} \\
Typst & \tscodeinline{N-\allowbreak{}01<n:joy>} & \tscodeinline{\#link(<n:joy>)[N-\allowbreak{}01]} \\
HTML & \tscodeinline{<span class="ts-\allowbreak{}counter" data-\allowbreak{}counter="n" data-\allowbreak{}key="joy" id="n:joy">N-\allowbreak{}01</span>} & \tscodeinline{<a href="\#n:joy">N-\allowbreak{}01</a>} \\
\bottomrule
\end{tabularx}
\end{center}
\tscodeinline{\textbackslash{}phantomsection} is what makes the \tscodeinline{hyperref} anchor land on the item rather
than on the enclosing section, and it lets \tscodeinline{\textbackslash{}pageref\{n:joy\}} work.
\section{On a MkDocs site}
The \tscodeinline{texsmith} MkDocs plugin renders counters on the site.
One plugin, no Markdown extension to wire:
\begin{tscode}[lang=yaml, engine=pygments]
\begin{Verbatim}[commandchars=\\\{\},breaklines, breakanywhere, commandchars=\\\{\}]
\PY{n+nt}{plugins}\PY{p}{:}
\PY{+w}{  }\PY{p+pIndicator}{\PYZhy{}}\PY{+w}{ }\PY{n+nt}{texsmith}\PY{p}{:}
\PY{+w}{      }\PY{n+nt}{declare}\PY{p}{:}
\PY{+w}{        }\PY{n+nt}{counters}\PY{p}{:}\PY{+w}{ }\PY{c+c1}{\PYZsh{} optional site\PYZhy{}wide declarations}
\PY{+w}{          }\PY{n+nt}{req}\PY{p}{:}
\PY{+w}{            }\PY{n+nt}{name}\PY{p}{:}\PY{+w}{ }\PY{l+lScalar+lScalarPlain}{Requirement}
\PY{+w}{            }\PY{n+nt}{format}\PY{p}{:}\PY{+w}{ }\PY{l+s}{\PYZdq{}}\PY{l+s}{REQ\PYZhy{}\PYZob{}n:03d\PYZcb{}}\PY{l+s}{\PYZdq{}}
\PY{+w}{            }\PY{n+nt}{start}\PY{p}{:}\PY{+w}{ }\PY{l+lScalar+lScalarPlain}{100}
\end{Verbatim}
\end{tscode}
Counters declared under \tscodeinline{declare.counters} apply to the whole site; a page may
declare its own under \tscodeinline{press.declare.counters} in its front matter, and the
page's declaration wins for the same prefix. Every \emph{item} is visible from every
page --- a series defined in \tscodeinline{findings.md} is referenceable from \tscodeinline{index.md}.

Numbering is \textbf{site-wide, in navigation order}: a pre-pass parses and resolves
every page before any of them is converted, chaining each page's counters after
the previous page's, so a reference on the first page resolves to an item
defined on the last one. Cross-page references are rewritten to point at the
page that defines the item (\tscodeinline{<a href="findings/\#fw:watchdog">FW-\allowbreak{}01</a>});
same-page ones keep a local anchor (\tscodeinline{<a href="\#fw:watchdog">FW-\allowbreak{}01</a>}).

The pre-pass parses the page rather than scanning it, so a marker inside a
fence or a code span is never counted, and \tscodeinline{\#\{user.name\}} in prose stays the
literal text it is.
\section{Citing an item from another document}
A number allocated here means nothing in a sister document, and a renumbering
breaks every hard-coded \tscodeinline{FW-\allowbreak{}10} it contains. Publish an inventory and cite it
explicitly --- see Cross-document references.
\section{Not implemented yet}
Auto-generated listings (a \enquote{List of requirements} table), extra per-item
attributes (status, priority, source), back-references ("cited on pages 4,
12"), hierarchical numbering (\tscodeinline{FW-\allowbreak{}3.2}) and per-chapter resets are deliberately
out of the first iteration.
\end{document}
